\section{A permanent expansion with a second-order uniform remainder}
\label{sec:secondpermanent}
The earlier \(N^{-1/8}\) remainder was sufficient for
\(A=o(\log N)\). At prime density it is necessary to retain
the first correction before summing the signed minors.

For a stochastic matrix \(D=J_N/N+B\), define
\[
 \Omega=\begin{pmatrix}0&B\\B^\top&0\end{pmatrix},\quad
 \Gamma=\Omega(I+\Omega)^{-1},\quad
 \Delta(s)=\det(I-s^2B^\top B)^{-1/2},
\]
and
\[
 \Psi(\Gamma)=
 -\frac34\sum_i\Gamma_{ii}^2
 +\frac12\sum_{i,j}\Gamma_{ii}\Gamma_{jj}\Gamma_{ij}
 +\frac13\sum_{i,j}\Gamma_{ij}^3,\qquad
 {\cal Q}_N(D)=\frac{\Delta''(1)}{2(N-1)\Delta(1)}
                                     +\Psi(\Gamma).
\]
The indices in \(\Psi\) range over the full \(2N\)-coordinate
block matrix, including all diagonal terms.

\begin{theorem}[First correction with a uniform two-power remainder]
\label{thm:secondpermanent}
Under Theorem~\ref{thm:growingpermanent}, for sufficiently large
\(N\) depending on \(\delta,K\),
\[
 \operatorname{per}D
 =\frac{N!}{N^N}\Delta(1)
   \left(1+{\cal Q}_N(D)
             +O_\delta((1+A)^{40}/N^2)\right),
 \qquad |{\cal Q}_N(D)|\leq C_\delta(1+A)^2/N .
\]
In particular the first correction is explicit and is not
hidden in a relative \(O(N^{-1/8})\) error.
\end{theorem}
\begin{proof}
Use the exact normalized Gaussian rook representation in
Section~\ref{sec:relativepermanents}. Its vertex factor has
expansion
\[
 r(z)=1+\tfrac13z^3-\tfrac14z^4+\tfrac15z^5
                              -\tfrac19z^6+O(z^7).
\]
Write \(\Gamma_s=s\Omega(I+s\Omega)^{-1}\).
The diagram with one degree-four vertex contributes
\(-3\sum_i(\Gamma_s)_{ii}^2/4\).
Two distinct degree-three vertices contribute
\[
 \sum_{i<j}
 \left((\Gamma_s)_{ii}(\Gamma_s)_{jj}(\Gamma_s)_{ij}
                              +\tfrac23(\Gamma_s)_{ij}^3\right).
\]
These coefficients follow from the three pairings of four
Gaussian variables and the \(9+6\) pairings of two triples.
Replacing \(i<j\) by half the ordered sum changes the result
by \(O_\delta(A^3/N^2)\), since
\(\max|(\Gamma_s)_{ii}|\leq C_\delta A/N\).
The one-vertex degree-six term also has that smaller order.

Here is a bound on all remaining diagrams, including their
infinite power-series tails.
In the finite Wick estimate already displayed, \(e\) is the
number of pairs and \(v\leq2e/3\).
For \(4\leq e\leq11\), one has \(e-v\geq2\);
their total is \(O_\delta((1+A)^{11}/N^2)\), after summing
the geometric coefficient convolution.
For \(e\geq12\), use the same cutoff
\(H=\lceil C_\delta A^3\sqrt N\rceil\) and sum
\[
 \sum_{e\geq12}e(C_\delta A^2N^{-1/6})^e
                    =O_\delta((1+A)^{24}/N^2).
\]
The cases \(e=1,2,3\) are exactly the terms retained above,
the zero odd Gaussian moments, and the smaller degree-six term.
The Cauchy tail past \(H\) is exponentially small.
This also applies to the retained degree-four and two-triple
contributions themselves: their covariance power series are
analytic on the same fixed circle \(|t|=R\), where the
resolvent bounds above bound their suprema by a polynomial
in \(1+A\). Hence extending their truncated coefficient sums
to their full values introduces only an exponentially small tail.
Thus, uniformly on the fixed interval \(s\in[a,b]\),
\[
 R_B(s)=\Delta(s)
        \left(1+\Psi(\Gamma_s)
                  +O_\delta((1+A)^{40}/N^2)\right).
\]
No coefficientwise asymptotic with an unspecified entry-bound
constant is used.

The row and column squared norms give
\[
 \|\Gamma_s\|\leq C_\delta,\quad
 \|\Gamma_s\|_F^2\leq C_\delta A,\quad
 \max_{ij}|(\Gamma_s)_{ij}|\leq C_\delta A/N.
\]
The same bounds hold for any fixed number of \(s\)-derivatives.
For \(g_i=(\Gamma_s)_{ii}\), the middle term in \(\Psi\) is
\(g^\top\Gamma_s g/2\), bounded by \(C_\delta A^2/N\).
The cubic entry sum is bounded by
\(\max|\Gamma_s|\|\Gamma_s\|_F^2\).
This proves \(|\Psi(\Gamma_s)|\leq C_\delta A^2/N\),
also for the derivatives needed below.

It remains to perform the Gamma average with its first
correction retained. Put \(S=N/T\), \(T\sim\Gamma(N+1,1)\),
\(I=[a,b]\), and
\(C_N=\{|S-1|\leq N^{-1/4}\}\subseteq\{S\in I\}\).
All expectations of \(\Delta(S)\) or
\(\Delta(S)\Psi(\Gamma_S)\) below are restricted to \(S\in I\).
No extension of these local analytic factors outside \(I\)
is needed.

For integers \(0\leq j\leq8\) and \(N\geq16\), direct
integration of the Gamma density gives
\[
 \mathbb E S^j=\prod_{\ell=0}^{j-1}\frac{N}{N-\ell}
 \leq2^j .
\]
In particular the exact full centered moments are
\[
 \mu_1=0,\qquad \mu_2=\frac1{N-1},\qquad
 \mu_3=\frac4{(N-1)(N-2)},\qquad
 \mu_4=\frac{3(N+6)}{(N-1)(N-2)(N-3)} .
\]
More explicitly, if \(0<\varepsilon\leq1/2\) and
\(N\varepsilon\geq4\), both actual Gamma tails give
\[
 \mathbb P(|S-1|>\varepsilon)
 \leq2\exp(-(N+1)\varepsilon^2/36).
\]
For \(4k\leq N\), the raw moments imply
\(\mathbb E|S-1|^{2k}\leq2^{2k}+1\).
Put \(u=\exp(-(N+1)\varepsilon^2/72)\) and
\(E_\varepsilon=\{|S-1|>\varepsilon\}\).
The elementary quadratic inequality
\[
 u|S-1|^k\one_{E_\varepsilon}
 \leq u^2|S-1|^{2k}+\one_{E_\varepsilon}
\]
therefore proves the explicit bound
\[
 \mathbb E[|S-1|^k\one_{E_\varepsilon}]
 \leq(2^{2k}+3)\exp(-(N+1)\varepsilon^2/72).
\]
Taking \(\varepsilon=N^{-1/4}\) gives the required
\(C_k\exp(-\sqrt N/72)\) bound for \(0\leq k\leq4\).
Thus each truncated signed moment on \(C_N\) equals
\(\mu_k+O(\exp(-c'\sqrt N))\), including the first
and third moments. In particular the truncated first moment
is not set to zero without its error.

For clarity, the precise local averaging estimate is as follows.
Let \(C_\varepsilon=\{|S-1|\leq\varepsilon\}\), \(M\geq0\),
and let \(f\) be measurable with
\[
 |f(s)-a-b(s-1)-c(s-1)^2-d(s-1)^3|
 \leq M(s-1)^4\quad (|s-1|\leq\varepsilon).
\]
For \(N\geq16\) and the preceding range of \(\varepsilon\),
the actual normalized Euler integral satisfies
\[
 \begin{split}
 &\left|\mathbb E[\one_{C_\varepsilon}f(S)]
       -a-\frac{c}{N-1}-\frac{4d}{(N-1)(N-2)}\right|\\
 &\quad\leq M\mu_4+
 e^{-(N+1)\varepsilon^2/72}
       (4|a|+7|b|+19|c|+67|d|).
 \end{split}
\]
Indeed the full cubic has exactly the displayed signed average;
only its discarded tail is bounded absolutely.
The local remainder is bounded by \(M\mu_4\).
These bounds also prove integrability on the window, without
extending the local approximation to other values of \(S\).

On \(C_N\), fixed-order derivatives of
\(\Delta(s)/\Delta(1)\) are bounded by polynomials in \(1+A\).
Indeed its logarithmic derivatives have that property and
\(\Delta(s)/\Delta(1)\leq\exp(C_\delta A N^{-1/4})\)
is bounded for sufficiently large \(N\).
Taylor's formula to third order, with a fourth-derivative
remainder and the truncated moments just proved, gives
\[
 \mathbb E\bigl[\one_{\{S\in I\}}\Delta(S)\bigr]
 =\Delta(1)\left(
  1+\frac{\Delta''(1)}{2(N-1)\Delta(1)}
                  +O_\delta((1+A)^8/N^2)\right).
\]
Here the cubic term uses the signed \(\mu_3=O(N^{-2})\);
the fourth-order remainder uses \(\mu_4=O(N^{-2})\).
The contribution of \(I\setminus C_N\) is at most
\(\Delta(1)\exp(C_\delta A-c'\sqrt N)\), which is
smaller than any fixed inverse power of \(N\) throughout
\(A\leq(\log N)^K\).
All implicit constants in these averages depend only on \(\delta\);
\(K\) enters only the sufficiently-large-\(N\) threshold.

For \(f(s)=\Delta(s)\Psi(\Gamma_s)\), the previous covariance
and derivative bounds give
\[
 \sup_{C_N}\frac{|f''(s)|}{\Delta(1)}
       \leq C_\delta(1+A)^6/N .
\]
Taylor's formula through the linear term, the exponentially
small truncated first moment, and
\(\mathbb E(S-1)^2=1/(N-1)\) therefore give
\[
 \mathbb E\bigl[\one_{\{S\in I\}}\Delta(S)\Psi(\Gamma_S)\bigr]
 =\Delta(1)\Psi(\Gamma)
       +O_\delta\bigl(\Delta(1)(1+A)^{12}/N^2\bigr).
\]
The portion in \(I\setminus C_N\) is again exponentially
small. The uniform rook-expansion remainder averages to
\(O_\delta(\Delta(1)(1+A)^{40}/N^2)\).
On \(S\notin I\) we retain the original polynomial \(R_B(S)\);
the preceding outside-interval estimate gives
\[
 \left|\mathbb E\bigl[\one_{\{S\notin I\}}R_B(S)\bigr]\right|
 \leq\exp(-c_\delta N+C_\delta\sqrt{NA}).
\]
Since \(\Delta(1)\geq1\), this too is negligible relative
to the claimed error. Finally the logarithmic derivative
formulas give
\(\Delta''(1)/\Delta(1)=O_\delta((1+A)^2)\).
Combining the restricted averages and the original polynomial
tail proves both conclusions.
\end{proof}
